\documentclass{article}
\usepackage[T1]{fontenc}
\usepackage{lmodern}
\usepackage{graphicx}
\usepackage{amsmath,amssymb}
\usepackage[a4paper,margin=2cm]{geometry}
\usepackage[absolute,overlay]{textpos}
\setlength{\TPHorizModule}{1mm}
\setlength{\TPVertModule}{1mm}
\usepackage[indonesian]{babel}
\usepackage{hyperref}
\setlength{\emergencystretch}{2em}
% unit: Halaman 88

\begin{document}

% === Halaman 88 (python/native) ===

88

• Dekripsi:

 - Mula-mula hitung m -1 yaitu 7–1 (mod 26) dengan memecahkan 7x  1 (mod 26)

   Solusinya: x  15 (mod 26) sebab 7  15 = 105 1(mod 26).

 -  Jadi, P  15 (C – 10) (mod 26) 

 c1 = 2    → p1  15  (2 – 10) = –120  10 (mod 26) 

(huruf ‘k’)

 c2 = 25  → p2  15  (25 – 10) = 225  17 (mod 26) (huruf ‘r’)

 c3 = 14  → p3  15  (14 – 10) = 60  8 (mod 26) 

(huruf ‘i’)

 c4 = 11  → p4  15  (11 – 10) = 15  15 (mod 26) 

(huruf ‘p’)

 c5 = 13  → p5  15  (13 – 10) = 45  19 (mod 26) 

(huruf ‘t’)

 c6 = 4    → p6  15  (4 – 10) = –90   14 (mod 26) 

(huruf ‘o’)


 Plainteks yang diungkap kembali:  kripto

\end{document}
